% Generated draft: requires mathematical and source review.
\documentclass[11pt]{article}
\usepackage[margin=25mm]{geometry}
\usepackage{fontspec}
\setmainfont{Latin Modern Roman}
\usepackage{amsmath,amssymb,mathtools,unicode-math}
\setmathfont{Latin Modern Math}
\usepackage{enumitem,needspace,xurl,hyperref}
\setlength{\parindent}{0pt}
\setlength{\parskip}{5pt}
\setlength{\emergencystretch}{3em}
\newcommand{\N}{\mathbb{N}}
\newcommand{\Z}{\mathbb{Z}}
\newcommand{\Q}{\mathbb{Q}}
\newcommand{\R}{\mathbb{R}}
\newcommand{\C}{\mathbb{C}}
\newcommand{\F}{\mathbb{F}}
\newcommand{\A}{\mathbb{A}}
\newcommand{\PP}{\mathbb{P}}
\newcommand{\E}{\mathbb{E}}
\newcommand{\eps}{\varepsilon}
\newcommand{\dd}{\,\mathrm{d}}
\begin{document}
\section*{UnsolvedMath ID: 2302002}
\textbf{Source problem code:} catalogue entry
\subsection{Definitions and mathematical statement}
For an entire function \(f\), let \(M(r,f)=\max_{|z|=r}|f(z)|\) and define its order by
\[
\rho(f)=\limsup_{r\to\infty}\frac{\log\log M(r,f)}{\log r}.
\]
Finite order means \(\rho(f)<\infty\). A path to infinity may be represented by a continuous curve \(\gamma:[0,\infty)\to\C\) with \(|\gamma(t)|\to\infty\). A prescribed asymptotic value \(w\) along that path means \(f(\gamma(t))\to w\).

The problem asks for a general construction, for a finite collection of prescribed paths \(C_1,\ldots,C_k\) and different values \(w_1,\ldots,w_k\), of an entire function satisfying
\[
\rho(f)<\infty,\qquad
\lim_{z\to\infty,\ z\in C_\nu}f(z)=w_\nu
\quad(1\le\nu\le k),
\]
and having minimal growth in the intended sense. For finite values the limits are ordinary complex limits; the excerpt does not specify whether \(\infty\) is also allowed as a value.

The admissible class of paths and compatibility assumptions are absent. In particular, paths carrying different limiting values cannot share points arbitrarily far out. No unrestricted existence claim for arbitrary intersecting or otherwise incompatible paths is imposed. ``Minimal growth'' is not defined: least order, type at a given order, and a sharp radial majorant are distinct optimization targets, and an infimum need not be attained. Thus a precise construction theorem requires the intended path hypotheses and growth objective.

The source says Kennedy's and Al-Katifi's known methods \emph{seem} to require all prescribed values equal, unless the paths are straight lines. This is preserved as the source's qualified comparison, not a verified limitation of every such method. The supplied OPEN-TRIAGE summary calls this an Arakelyan/asymptotic-path construction problem, reports it open in the 2018 edition, and says no general resolution was located as of August 2026.
\subsection{Short English statement}
Construct an entire function of finite, optimally small growth with different prescribed limits along prescribed paths to infinity. The allowable paths, compatibility conditions, and meaning of minimal growth are not specified.
\subsection{Sources}
\begin{itemize}
\item Local UnsolvedMath registry, permanent ID \(2302002\), exact source code \texttt{catalogue entry}, ``Research Problems in Function Theory --- Problem 2.2.''
\item Hayman, \emph{Research Problems in Function Theory} (2018), Problem 2.2; \url{https://arxiv.org/abs/1809.07200}. Source-TeX extraction; recorded access 2026-07-29.
\item The source names methods by Kennedy and Al-Katifi without complete citations. The supplied summary describes an Arakelyan-type problem and an August 2026 search without a general resolution, but supplies no additional precise references.
\end{itemize}
\end{document}
